% ==================================================
% components_spec.tex
% L3 COMPONENT LAYER (spec page)
% ==================================================

\usepackage{array}
\usepackage{tabularx}
\usepackage{colortbl}
\usepackage[most]{tcolorbox}

% ---- fallbacks (safe defaults if params are missing) ----
\expandafter\providecommand\csname HBcomp_spec_section_before\endcsname{0.6mm}
\expandafter\providecommand\csname HBcomp_spec_first_section_before\endcsname{0.6mm}
\expandafter\providecommand\csname HBcomp_spec_section_after\endcsname{0.9mm}
\expandafter\providecommand\csname HBcomp_spec_section_bullet_raise\endcsname{-0.16ex}
\expandafter\providecommand\csname HBcomp_spec_section_bullet_symbol\endcsname{\scalebox{1.67}{\textbullet}}
\expandafter\providecommand\csname HBcomp_spec_section_bullet_gap\endcsname{0.36em}
\expandafter\providecommand\csname HBcomp_spec_table_left_ratio\endcsname{0.33}
\expandafter\providecommand\csname HBcomp_spec_table_tabcolsep\endcsname{3.2pt}
\expandafter\providecommand\csname HBcomp_spec_table_row_stretch\endcsname{1.10}
\expandafter\providecommand\csname HBcomp_spec_table_tall_row_extra\endcsname{2.1pt}
\expandafter\providecommand\csname HBcomp_spec_table_multiline_min_height\endcsname{23.4pt}
\expandafter\providecommand\csname HBcomp_spec_table_multiline_depth\endcsname{7.0pt}
\expandafter\providecommand\csname HBcomp_spec_environment_extra_before\endcsname{1.2mm}
\expandafter\providecommand\csname HBcomp_spec_notes_before\endcsname{0.8mm}
\expandafter\providecommand\csname HBcomp_spec_footnotes_before\endcsname{1.2mm}

\expandafter\providecommand\csname HBtype_spec_section_font_size\endcsname{8.8pt}
\expandafter\providecommand\csname HBtype_spec_section_font_leading\endcsname{9.6pt}
\expandafter\providecommand\csname HBtype_spec_label_font_size\endcsname{6.2pt}
\expandafter\providecommand\csname HBtype_spec_label_font_leading\endcsname{7.0pt}
\expandafter\providecommand\csname HBtype_spec_value_font_size\endcsname{6.2pt}
\expandafter\providecommand\csname HBtype_spec_value_font_leading\endcsname{7.0pt}
\expandafter\providecommand\csname HBtype_spec_note_font_size\endcsname{5.9pt}
\expandafter\providecommand\csname HBtype_spec_note_font_leading\endcsname{6.7pt}

% ---- spec section title (shared level-2 dot heading) ----
\newif\ifHBSpecFirstSection
\newcount\HBSpecSectionIndex
\providecommand{\specsectiontitle}[1]{%
  \par
  \global\advance\HBSpecSectionIndex by 1\relax
  \ifHBSpecFirstSection
    \vspace{\csname HBcomp_spec_first_section_before\endcsname}%
    \HBSpecFirstSectionfalse
  \else
    \vspace{\csname HBcomp_spec_section_before\endcsname}%
    \ifnum\HBSpecSectionIndex=4\relax
      \vspace{\csname HBcomp_spec_environment_extra_before\endcsname}%
    \fi
  \fi
  \Needspace{\csname HBcomp_title_l2_needspace\endcsname}%
  \HBTitleLevelTwo{#1}%
  \vspace{\csname HBcomp_spec_section_after\endcsname}%
}

% Specifications are a fixed-format page in the reference manual. These
% boundaries keep adjacent sections from flowing into or out of the page.
\providecommand{\HBSpecPageStart}{%
  \clearpage
  \global\HBSpecSectionIndex=0\relax
  \HBSpecFirstSectiontrue
}
\providecommand{\HBSpecPageEnd}{\clearpage}

% ---- spec table wrapper ----
% Each section fits on the dedicated spec page, so a non-breaking framed
% table can own the rounded outer stroke while tabularx owns only the inner
% divider and row rules.
\NewEnviron{spectable}{%
  \begingroup
  \parindent=0pt\relax
  \arrayrulecolor{HBSharedTableFrameColor}
  \let\tabularxcolumn\HBsharedXcolumn
  \setlength{\tabcolsep}{\csname HBcomp_spec_table_tabcolsep\endcsname}
  \renewcommand{\arraystretch}{\csname HBcomp_spec_table_row_stretch\endcsname}
  \arrayrulewidth=\dimexpr\csname HBcomp_table_outer_rule\endcsname/2\relax
  \ifdim\arrayrulewidth<0.24pt\arrayrulewidth=0.24pt\fi
  \begin{tcolorbox}[
    enhanced,
    clip upper,
    colback=white,
    colframe=HBSharedTableFrameColor,
    arc=\csname HBcomp_table_outer_arc\endcsname,
    boxrule=\csname HBcomp_table_outer_rule\endcsname,
    boxsep=0mm,
    left=0mm,right=0mm,top=0mm,bottom=0mm,
    before skip=0mm,after skip=0mm
  ]%
  \begin{tabularx}{\linewidth}{@{}
    >{\columncolor{HBSharedTableLeftBg}\RaggedRight\arraybackslash}%
      m{\csname HBcomp_spec_table_left_ratio\endcsname\linewidth}|
    >{\RaggedRight\arraybackslash}X@{}}%
    \BODY
  \end{tabularx}%
  \end{tcolorbox}%
  \endgroup
}

% Reference rows containing multiple value lines or a footnote marker carry
% a little more vertical air than ordinary single-line rows.
\providecommand{\HBSpecTallRowBreak}{%
  \tabularnewline[\csname HBcomp_spec_table_tall_row_extra\endcsname]%
}
\providecommand{\HBSpecMultilineRowStrut}{%
  \rule[-\csname HBcomp_spec_table_multiline_depth\endcsname]{0pt}%
       {\csname HBcomp_spec_table_multiline_min_height\endcsname}%
}

% ---- spec text styles ----
\providecommand{\HBTypeSpecLabel}[1]{%
  {\normalfont\mdseries
   \color{\csname HBtype_body_color\endcsname}%
   \fontsize{\csname HBtype_spec_label_font_size\endcsname}
            {\csname HBtype_spec_label_font_leading\endcsname}
   \selectfont
   \hspace*{\csname HBcomp_table_outer_arc\endcsname}#1}%
}

\providecommand{\HBTypeSpecValue}[1]{%
  {\normalfont\mdseries
   \color{\csname HBtype_body_color\endcsname}%
   \fontsize{\csname HBtype_spec_value_font_size\endcsname}
            {\csname HBtype_spec_value_font_leading\endcsname}
   \selectfont #1}%
}

\providecommand{\HBTypeSpecNote}[1]{%
  {\normalfont\mdseries
   \color{\csname HBtype_body_color\endcsname}%
   \fontsize{\csname HBtype_spec_note_font_size\endcsname}
            {\csname HBtype_spec_note_font_leading\endcsname}
   \selectfont #1}%
}

% ---- portable note/footnote markers ----
% Avoid tiny math fallback fonts from \textsuperscript in some XeLaTeX stacks.
\providecommand{\HBSpecMarkerBase}[1]{%
  \raisebox{0.48ex}{\fontsize{5.2pt}{5.2pt}\selectfont #1}%
}
\providecommand{\HBSpecMarkerOne}{\HBSpecMarkerBase{1}}
\providecommand{\HBSpecMarkerTwo}{\HBSpecMarkerBase{2}}
\providecommand{\HBSpecMarkerThree}{\HBSpecMarkerBase{3}}
\providecommand{\HBSpecMarkerFour}{\HBSpecMarkerBase{4}}
\providecommand{\HBSpecMarkerFive}{\HBSpecMarkerBase{5}}
\providecommand{\HBSpecMarkerSix}{\HBSpecMarkerBase{6}}
\providecommand{\HBSpecMarkerSeven}{\HBSpecMarkerBase{7}}
\providecommand{\HBSpecMarkerEight}{\HBSpecMarkerBase{8}}
\providecommand{\HBSpecMarkerNine}{\HBSpecMarkerBase{9}}
\providecommand{\HBSpecMarkerTen}{\HBSpecMarkerBase{10}}
\providecommand{\HBSpecMarkerAsterisk}{\HBSpecMarkerBase{*}}
